

\SourceLeaf{221}%
Weil in Bezug auf diesen Halbkreis $a$, $b$, $c$ als verschwindend zu betrachten sind, also die Punkte $A$, $B$ als in $C$ fallend, so ist dieser Halbkreis das Mass der drei Winkel des Dreiecks, die mithin weniger als jede gegebene Grösse von $180^0$ differiren.

Mir scheint, wenn man den Begriff des endlos Wachsenden nicht ausschliesst, so zeigt dieser Beweis sehr einfach, dass in jedem \textls{endlichen} geradlinigen Dreiecke die Summe der Winkel $=180^0$ ist, oder eigentlich, dass die Constante, die, wenn Euklids Geometrie nicht wahr wäre, zu der Summe der Winkel kommt, um die Gleichheit mit $180^0$ zu bewirken, kleiner als jede gegebene Grösse ist, und da sich diess für jedes Dreieck beweisen lässt, so kann diese Constante eben so wenig von der Grösse des Dreiecks abhängen. \mbox{.\ .\ .\ .\ .\ .}\}

\bigskip
\hrule
\bigskip

\Needspace{7\baselineskip}
\begin{center}
[\textsc{Schumacher} \emph{an} \textsc{Gauss}. \emph{Altona, 29. Juni 1831.}]
\end{center}

\{\mbox{.\ .\ .\ .\ .\ .}\ .\ .\ Nur etwas habe ich in Ihrem Briefe vermisst — Ihr Urtheil über meinen Beweis, dass die Summe der Winkel in einem geradlinigen Dreiecke nur um eine Grösse, die kleiner als jede gegebene ist, von $180^0$ verschieden sei. Sie können leicht denken, dass mir Ihr Urtheil sehr wichtig ist, da Sie jede Schwäche eines Beweises so leicht entdecken. Ausser Ihnen, meinem Gehülfen, und Professor \textsc{Hansen} vom Seeberg habe ich noch niemandem etwas mitgetheilt. Keiner von uns kann einen Paralogismus entdecken.

Sollte jemand den Satz, dass man die Winkelpunkte eines endlichen Dreiecks als coincidirende Mittelpunkte eines Kreises von unendlichem (brevitatis causa \textls{unendlich} genannt) Halbmesser betrachten könne, eines Beweises bedürfend halten, obgleich ich diess nicht glaube, so lässt sich dieser Beweis strenge führen.

Mir scheint, wenn zwei Punkte eine endliche Entfernung von einander haben, so wird diese Entfernung in Bezug auf eine unendliche Linie $=0$ zu setzen sein, sie coincidiren mithin in Bezug auf diese unendliche Linie betrachtet. \mbox{.\ .\ .\ .\ .\ .}\ .\}

\bigskip
\hrule
\bigskip
